\section{A derivative algorithm with fixed schedules}
\label{sec:fixed-schedule-derivative}

We fix the computation first and prove its validity separately.  A stage
must not mean ``keep refining until an error test succeeds.''  Choose
explicit rational increments $h_n(x)>0$ tending to zero and explicit
natural-number stages $\sigma(n,x)$, with $x\pm h_n(x)$ in the domain.
Typical choices are $h_n=h_0/(n+1)$ and $\sigma(n,x)=(n+1)^2$; a
particular evaluator may permit the simpler $\sigma(n,x)=n$.

The notation $[E]_n$ means the two rational outputs of a specified
composite expression at its prescribed stage.  It does not denote a scalar
$E$, and it does not silently introduce a precision search.  For example,
using ordinary stagewise subtraction,
\[
 [f(x)-f(y)]_n=
 [f_n^-(x)-f_n^+(y),\ f_n^+(x)-f_n^-(y)].
\]
If a composite evaluator uses a different internal schedule, that schedule
is stated as part of its definition.

\begin{definition}[Fixed-schedule derivative evaluator]
\label{def:computed-derivative}
For a convex evaluator $f$, set, with $h=h_n(x)$ and $s=\sigma(n,x)$,
\[
 B_n^{h,\sigma}(x)=
 \left[\frac{[f(x)-f(x-h)]_s^-}{h},
       \frac{[f(x+h)-f(x)]_s^+}{h}\right],
\]
and compute the finite intersections
\[
 [D_{h,\sigma}f(x)]_n=
 \left[\max_{j\le n}(B_j^{h,\sigma}(x))^-,
       \min_{j\le n}(B_j^{h,\sigma}(x))^+\right].
\]
All loops have a length determined by $n$; there is no error-dependent
stopping condition.  This defines an evaluator before asserting that its
widths shrink.
\end{definition}

Endpoint convexity proves $(B_i^{h,\sigma})^-\le(B_j^{h,\sigma})^+$
for every $i,j$.  Indeed a left secant at $x$ is below every right secant,
by the three-point rational convexity inequality.  Coarse lower endpoints
can only decrease the former, and coarse upper endpoints can only increase
the latter.  Refining the finitely many evaluations in this proof and
removing rational slack establishes the inequality at the originally fixed
stages.  Thus every finite intersection is ordered, and nesting follows
from finite maxima and minima.  Neither a derivative scalar nor an infinite
intersection witness is used.

\begin{theorem}[Validity for a specified schedule]
\label{thm:fixed-schedule-derivative}
Suppose the rational widths
$\operatorname{width}B_n^{h,\sigma}(x)$ tend to zero for the specified
convex evaluator and schedules.  Then $D_{h,\sigma}f$ is valid at $x$
and differentiates $f$ in the endpoint-residual sense.  Any two successful
schedules, including schedules used with equivalent source evaluators,
produce derivative pairs overlapping at every pair of stages.
\end{theorem}
\begin{proof}
The finite intersection has width no greater than the current bracket.
For the derivative assertion fix a positive rational slope tolerance.
Choose one outer bracket narrower than a fixed fraction of that tolerance.
Convexity brackets every nonzero rational increment shorter than its $h_j$
by the same outer secants.  A sufficiently late derivative output lies in
that chosen bracket.  For the particular increment under consideration,
refine the finitely many source outputs in the proof until their widths
are smaller than the tolerance times the increment's absolute value.
The resulting residual pair lies inside the prescribed rational error
interval.  This establishes the derivative relation, whose witnessing
stage may depend on the increment; it does not change the evaluator's
fixed execution rule.  The endpoint proof of derivative uniqueness then
gives the final assertion.
\end{proof}

Orderedness and termination of each finite stage must be distinguished from
validity.  For the exact evaluator of $|x|$ at zero all these brackets are
$[-1,1]$.  Every stage runs, but no shrinking-width theorem is true there.

\subsection{Resolution of the source and size of the increment}
Write $w_m(t)=f_m^+(t)-f_m^-(t)$, a rational number.  The relevant
resolution budget is
\[
 \rho_n(x)=\frac{
 w_{\sigma(n,x)}(x-h_n)+2w_{\sigma(n,x)}(x)
                 +w_{\sigma(n,x)}(x+h_n)}{h_n}.
\]
Smaller increments amplify source uncertainty.  Validity of $f$ at each
fixed input alone gives no rate at the moving inputs $x\pm h_n$.

If an independently specified valid evaluator $g$ differentiates $f$,
and $\rho_n(x)\to0$, then the brackets above shrink.  To prove this,
apply the two residual estimates at $h_n$ and $-h_n$, at sufficiently
refined proof-side stages.  Their secant gap is bounded by twice the
slope tolerance and a derivative-output width.  Returning to the prescribed
source stage adds at most $\rho_n(x)$.  Remove the arbitrarily small
proof-side width and then reduce the tolerance.  Unlike an adaptive
construction, an arbitrary fixed schedule does \emph{not} inherit
convergence from differentiability without a resolution condition.

For a concrete smooth computation one often proves
\[
 \operatorname{width}B_n^{h,\sigma}(x)
       \le K h_n+\rho_n(x)
\]
by a finite second-difference identity.  The constant $K$ and the source
width bound are theorems about that computation, not runtime inputs
requesting an accuracy test.

\begin{example}[Two schedules for the same noisy square]
For $m\ge1$, take $f_m(t)=[t^2-m^{-2},t^2+m^{-2}]$ and use
$\sigma(n)=n$.  Exact rational arithmetic gives
\[
 B_n^h(x)=
 [2x-h_n-2/(n^2h_n),\ 2x+h_n+2/(n^2h_n)].
\]
For $h_n=1/n$ the width is $6/n$; for $h_n=1/n^2$ it is $4+2/n^2$.
The latter brackets are already nested, so finite intersections do not
repair their failure to shrink.  With the exact square evaluator both
choices work.  The validity theorem therefore concerns the evaluator
\emph{together with} its schedules.
\end{example}

A uniform source-width bound $w_m(t)\le C/(m+1)$ on the sampling interval
can be paired with $h_n=h_0/(n+1)$ and
$\sigma(n)=(n+1)^2$: the resolution budget is at most
$4C/(h_0(n+1))$.  Higher polynomial schedules give correspondingly smaller
errors.  These are predetermined integer formulas, not precision searches.

\subsection{Correcting curvature without changing the principle}
If the finite inequalities prove that $f+Kx^2$ is convex, use its secants
and subtract the exact derivative $2Kx$.  With exact polynomial arithmetic
the resulting bracket can be written directly as
\[
 \left[
 \frac{[f(x)-f(x-h_n)]_{\sigma(n)}^-}{h_n}-Kh_n,
 \quad
 \frac{[f(x+h_n)-f(x)]_{\sigma(n)}^+}{h_n}+Kh_n
 \right].
\]
This is convenient for sine on an interval crossing a change of curvature.
The polynomial terms cancel before interval evaluation; separately evaluating
equal occurrences would unnecessarily enlarge the pairs.

\subsection{A closed factorial-series example}
\label{sec:fixed-factorial-sine}
The following is an alternative representation, constructed from finite
polynomials rather than from inverse sector area.  Put $N=n+5$ and
\[
 s_n(t)=\sum_{j=0}^{N-1}\frac{(-1)^jt^{2j+1}}{(2j+1)!},\qquad
 c_n(t)=\sum_{j=0}^{N-1}\frac{(-1)^jt^{2j}}{(2j)!},\qquad
 r_n=4\frac{2^{2N}}{(2N)!}.
\]
These are rational numbers for rational $t$.  On $[-2,2]$ the pairs
$[s_n(t)-r_n,s_n(t)+r_n]$ and $[c_n(t)-r_n,c_n(t)+r_n]$ are nested
and shrinking.  The factorial-tail comparison proves nesting, not merely
agreement with an externally supplied sine or cosine.  The series chapter
gives the general finite-tail argument.

For $h_n=1/(n+1)$ one has $r_n\le h_n^2$.  Indeed the tail-term ratio
from index $10$ onward is at most $1/2$, so
\[
 r_n\le4\frac{2^{10}}{10!}\,2^{-2n}\le(n+1)^{-2}.
\]
The last inequality follows from $2^n\ge n+1$ by induction.

Finite polynomial calculations on $|t|,|t+h|\le2$ give
\[
 \left|\frac{s_n(t+h)-s_n(t)}h-c_n(t)\right|\le34|h|,
 \qquad |c_n(t)-c_n(u)|\le16|t-u|.
\]
For an elementary derivation, $s_n'=c_n$ is a coefficient identity.
The absolute derivative polynomials are bounded by the finite exponential
majorant at $2$, which is less than $8$: terms through degree four sum
to $7$, and the remaining terms sum to at most $2/5$ by their ratio
bound.  The polynomial FTC and its twice-integrated remainder give even
smaller constants; the displayed conservative constants suffice here.
This invokes only the previously proved finite polynomial identities.

Consequently, for $|x|\le1$, each left and right center secant is within
$50h_n$ of $c_n(x)$.  Define
\[
 B_n(x)=\left[
 \frac{s_n(x)-s_n(x-h_n)-2r_n}{h_n}-50h_n,
 \quad
 \frac{s_n(x+h_n)-s_n(x)+2r_n}{h_n}+50h_n
 \right].
\]
It uses only the sine pair evaluator, the fixed quadratic correction,
and the same stage $n$.  No cosine evaluation occurs in its definition.
The finite secant estimates prove that $B_n(x)$ contains the entire
cosine pair $[c_n(x)-r_n,c_n(x)+r_n]$, since $h_n\le1$ and
$2r_n/h_n\ge r_n$.  They also give
\[
 \operatorname{width}B_n(x)\le200h_n+4r_n/h_n\le204/(n+1).
\]
The prefix intersections contain the stage-$n$ cosine pair: later cosine
pairs lie in earlier ones, so each is in all earlier $B_j$.  Thus the
intersections are valid and overlap the cosine evaluator at every pair
of stages.  The all-increment polynomial estimate and the factorial tails
also give the endpoint derivative relation, not just convergence along
one sequence of increments.

This proves the derivative comparison for the radian-coordinate factorial
representation.  In the book's normalized angle coordinate the intended
identity is $D S=\pi C$.  Multiplication by the computed $\pi$ and
evaluation on its shrinking rational input pairs require their own fixed
schedules and domain bounds; identifying that representation with the
geometric circle construction is a separate comparison.  A radian-series
proof must not be advertised as completing that geometric comparison.

\subsection{Using a fixed derivative in an integral}
An integral stage may likewise prescribe a finite partition and the
stages used for its derivative samples.  Its proof then bounds the mesh
error and the finite sum of sample widths.  For example, a derivative
width bound $C/(m+1)$ permits polynomial reindexing of those samples.
No stopping test is required inside the derivative evaluator.  The
endpoint-only FTC below identifies the resulting integral whenever the
chosen rectangle schedule has a proved shrinking-width bound.
